%% ---------------------------------------------------------------------
\section{Estabilidad: de local a global}
\label{sec:fl-estabilidad}
%% ---------------------------------------------------------------------

\subsection{Fase mínima basta, localmente}

\begin{teorema}[label=thm:fl-local]{Estabilización local \citep{ByIs88}}
Suponga que el sistema \eqref{eq:fl-afin} tiene grado relativo $r$ en el equilibrio $x^o$, con $h(x^o)=0$, y que $\eta^o$ es un equilibrio \emph{localmente exponencialmente estable} de la dinámica cero \eqref{eq:fl-cero}. Entonces la ley
\[
  u=\frac{1}{b(z)}\bigl(-a(z)+v\bigr),\qquad v=-k_1\xi_1-k_2\xi_2-\dots-k_r\xi_r,
\]
con $s^r+k_rs^{r-1}+\dots+k_1$ Hurwitz, hace $x^o$ localmente exponencialmente estable.
\end{teorema}

\begin{proof}
En coordenadas $(\xi,\eta)$, el lazo cerrado es $\dot\xi=(A_c-B_cK)\xi$, $\dot\eta=q(\xi,\eta)$. Su jacobiano en $(0,\eta^o)$ es triangular por bloques,
\[
  J=\begin{bmatrix}A_c-B_cK&0\\[2pt] \dfrac{\partial q}{\partial\xi}(0,\eta^o)&\dfrac{\partial q}{\partial\eta}(0,\eta^o)\end{bmatrix},
\]
y sus valores propios son los de $A_c-B_cK$ (Hurwitz por diseño) y los de $\frac{\partial q}{\partial\eta}(0,\eta^o)$ (Hurwitz, porque la estabilidad exponencial de la dinámica cero equivale a que su linealización sea Hurwitz). Por el método indirecto de Lyapunov, $(0,\eta^o)$ es localmente exponencialmente estable, y como $\Phi$ es un difeomorfismo local, también lo es $x^o$.
\end{proof}

\begin{cuidado}[Asintótica no es exponencial]
Si la dinámica cero es asintóticamente estable pero no exponencialmente (por ejemplo $\dot\eta=-\eta^3$), $\frac{\partial q}{\partial\eta}(0,\eta^o)$ tiene un valor propio en el eje imaginario, y el método indirecto no concluye nada. La estabilidad local se puede probar todavía, pero con una función de Lyapunov (o con el teorema de la variedad central); no con la jacobiana.
\end{cuidado}

\begin{ejemplo}[label=ej:fl-faseminima]{Un sistema de fase mínima}
Estabilice el origen del sistema \eqref{eq:fl-ejkhalil}.
\solucion Con la forma normal del \cref{ej:fl-formakhalil}, anular la salida es $z_1=z_2=0$, es decir
\[
  Z^*=\{x\in\R^3:\ x_1=0,\ x_3=x_2^3\},
\]
y la dinámica cero es $\dot z_3=-z_3$: globalmente exponencialmente estable. El sistema es de fase mínima. Como $b=1+3x_2^2$ nunca se anula, la ley
\[
  u=\frac{v-(3x_2^3+x_1^2-x_3)}{1+3x_2^2},\qquad v=-2x_1-3(x_3-x_2^3),
\]
(polos en $-1$ y $-2$) está definida en todo $\R^3$, y por el \cref{thm:fl-local} el origen es localmente exponencialmente estable. En $Z^*$, $z_3=x_2+x_3\to0$ y $x_3=x_2^3$, así que $x_2,x_3\to0$. Desde $x(0)=(1,\ 0.5,\ -0.5)$, $\norm{x(10)}<10^{-3}$.
\end{ejemplo}

\codigorepo{cap12_linealizacion/C12_Ej11_DinamicaCero.m}{faseminima}{Un sistema de fase mínima con la ley linealizante}

\subsection{Global pide ISS}

El \cref{thm:fl-local} es local, aunque la dinámica cero sea \emph{globalmente} exponencialmente estable. Para una conclusión global hay que controlar cómo $\xi$ empuja a $\eta$ mientras converge: el lazo cerrado es una cascada, y la herramienta es la estabilidad entrada--estado.

\begin{proposicion}[label=prop:fl-global]{Estabilización global \citep[Lema~4.7]{Kha02}}
Si en coordenadas globales el subsistema $\dot\eta=q(\xi,\eta)$ es ISS con $\xi$ como entrada, la ley del \cref{thm:fl-local} hace el origen globalmente asintóticamente estable.
\end{proposicion}

La prueba es el resultado de cascadas de la unidad de ISS: $\xi(t)\to0$ exponencialmente, y un sistema ISS con una entrada que tiende a cero tiene su estado acotado y convergente a cero. Sin ISS, la conclusión falla aun en sistemas muy sencillos.

\begin{trampa}[label=tr:fl-global]{``Dinámica cero globalmente estable, estabilización global''}
\plausible{En $\dot\xi=-k\xi$, $\dot\eta=-\eta+\eta^2\xi$ la dinámica cero es $\dot\eta=-\eta$ (globalmente exponencialmente estable) y $\xi\to0$ exponencialmente; luego todo converge.}
El origen es localmente exponencialmente estable (\cref{thm:fl-local}), pero no global. Con $w=1/\eta$, $\dot w=w-\xi_0e^{-kt}$, y
\[
  w(t)=e^{t}\Bigl(w_0-\frac{\xi_0}{k+1}\Bigr)+\frac{\xi_0}{k+1}\,e^{-kt} .
\]
Si $\eta_0\xi_0>k+1$, $w$ se anula en tiempo finito: $\eta$ \emph{escapa a infinito}, en $t_e=\frac{1}{k+1}\ln\dfrac{\xi_0/(k+1)}{\xi_0/(k+1)-1/\eta_0}$. Con $k=1$ y $\xi_0=\eta_0=2$, $t_e=\tfrac12\ln2=0.3466$. Subir $k$ agranda la región que converge, pero nunca la vuelve global (\cref{fig:fl-escape}). Es la misma cascada de la \cref{tr:pa-cascada}: el subsistema $\eta$ no es ISS respecto a $\xi$, porque el término $\eta^2\xi$ crece más rápido de lo que $-\eta$ puede contener.
\correcto{La estabilidad global de la dinámica cero da estabilización local. Para la global hace falta ISS de $\dot\eta=q(\xi,\eta)$ respecto a $\xi$ (\cref{prop:fl-global}), o una ley que tenga en cuenta el acoplamiento (backstepping, pasivación: la \cref{thm:pa-cascada}).}
\end{trampa}

\begin{figure}[ht]
\centering
\includegraphics[width=\linewidth]{cap12/f_escape}
\caption{Trayectorias de $\dot\xi=-k\xi$, $\dot\eta=-\eta+\eta^2\xi$ en el plano $(\eta,\xi)$. Las que empiezan en la región sombreada ($\eta_0\xi_0>k+1$) escapan en tiempo finito (rojas). Con $k=5$ la región es más pequeña, pero no desaparece.}
\label{fig:fl-escape}
\end{figure}

\codigorepo{cap12_linealizacion/C12_Ej12_Global.m}{iss}{Escape en tiempo finito con dinámica cero estable}

\subsection{No cancelar lo que ayuda}

\begin{trampa}[label=tr:fl-cancelar]{``Linealizar por retroalimentación es cancelar todas las no linealidades''}
\plausible{Para $\dot x=x-x^3+u$, la ley natural es cancelar: $u=x^3-2x$ da $\dot x=-x$, un sistema lineal y estable.}
El término $-x^3$ \emph{ya} estabiliza cuando $\abs x$ es grande. Cancelarlo es pagar por quitar algo que ayudaba. Compare con $u=-2x$, que da $\dot x=-x-x^3$ (también GAS, con $V=\tfrac12x^2$). Desde $x(0)=3$: la cancelación pide $u(0)=21$ (contra $-6$); gasta $\int_0^\infty u^2\,\dif t=117/2=58.5$ (contra $2\ln10=4.6$, unas $12.7$ veces más), y converge más lento: $x(1)=3/e=1.104$ contra $0.372$ (\cref{fig:fl-nocancelar}).
\correcto{La linealización por retroalimentación es una herramienta de análisis y de diseño, no una obligación de cancelar todo. Los diseños ``constructivos'' (funciones de Lyapunov de control, backstepping, pasivación) conservan las no linealidades útiles~\citep{KoAr01}.}
\end{trampa}

\begin{figure}[ht]
\centering
\includegraphics[width=\linewidth]{cap12/f_no_cancelar}
\caption{$\dot x=x-x^3+u$ desde $x(0)=3$: cancelar (rojo) contra conservar el término estabilizante (azul). Estado, control y energía de control acumulada.}
\label{fig:fl-nocancelar}
\end{figure}

\codigorepo{cap12_linealizacion/C12_Ej12_Global.m}{nocancelar}{Cancelar contra conservar}

\subsection{El puente con la pasividad}

El \cref{cap:pasividad} terminó con una pregunta: ¿cuándo se puede volver pasivo un sistema con una realimentación preliminar? La respuesta está escrita en el lenguaje de este capítulo.

\begin{teorema}[label=thm:fl-biw]{Pasividad y fase mínima \citep{BIW91}}
El sistema \eqref{eq:fl-afin} es equivalente por retroalimentación (con $u=\alpha(x)+\beta(x)v$, $\beta\neq0$) a un sistema pasivo con almacenamiento definido positivo si y solo si tiene grado relativo uno en $x=0$ y es de fase mínima débil.
\end{teorema}

Las dos condiciones tienen sentido a la luz del \cref{cap:pasividad}. Si $\Vdot\le u\,y$ con $y=h=L_gV$, la entrada aparece en $\dot y$ a través de $L_gh=L_g(L_gV)$, que en el origen vale $g\T\,\nabla^2V\,g$: positiva si el hessiano de $V$ es definido positivo y $g(0)\neq0$, de modo que $r=1$. Y en $Z^*=\{y=0\}$, $\Vdot\le0$: la dinámica cero tiene a $V$ como función de Lyapunov, y es al menos débilmente estable. Con la equivalencia, la ley $v=-\phi(y)$ estabiliza por inyección de amortiguamiento (\cref{thm:pa-amortiguamiento}). Las dos últimas unidades del curso son dos caras de la misma estructura: la \emph{energía} (pasividad) y la \emph{geometría} (grado relativo y dinámica cero).

\begin{profundiza}[¿Cuándo existe una salida con $r=n$?]
Si existe una salida $h$ con grado relativo $n$, el sistema es linealizable entrada--estado y no hay dinámica cero. Jakubczyk y Respondek~\citep{JaRe80} y Hunt, Su y Meyer~\citep{HSM83} probaron que esa $h$ existe cerca de $x^o$ si y solo si
\begin{enumerate}
  \item la matriz $\bigl[\,g\ \ \mathrm{ad}_fg\ \ \cdots\ \ \mathrm{ad}_f^{n-1}g\,\bigr](x^o)$ tiene rango $n$ (el análogo de la controlabilidad), y
  \item la distribución $\operatorname{span}\{g,\mathrm{ad}_fg,\dots,\mathrm{ad}_f^{n-2}g\}$ es \emph{involutiva} cerca de $x^o$ (cerrada bajo el corchete de Lie),
\end{enumerate}
con $\mathrm{ad}_fg=[f,g]=\jac gxf-\jac fxg$ \citep[\S13.2]{Kha02}, \citep[\S4.2]{Isi95}. Aun si existe, $h$ puede no tener sentido físico: la salida que linealiza no tiene por qué ser la que se quiere regular. (Fuera del alcance de la evaluación del curso.)
\end{profundiza}
